\subsubsection{Route-level iterative proportional fitting baseline}
\label{sec:reduced-od-route-level-ipf}

The first numerical reduced-OD baseline estimates vehicle-leg OD flows for one
route pattern and service class.  It is intentionally simpler than the network
journey model and is implemented in
\path{public_transportation.inference.reduced_od.route_level}.  Consider an
ordered route with stops $1,\ldots,n$.  Let $x_{ij}$ be the passenger flow that
boards at stop $i$ and alights at stop $j$.  The structural support is strictly
upper triangular:
\[
x_{ij}\geq 0\quad\text{for }i<j,
\qquad
x_{ij}=0\quad\text{for }i\geq j.
\]
For observed boarding marginals $b_i$ and alighting marginals $a_j$, iterative
proportional fitting alternately scales constrained rows and columns toward
\[
\sum_{j>i}x_{ij}=b_i,
\qquad
\sum_{i<j}x_{ij}=a_j.
\]
A supplied seed must be finite and strictly positive on every structurally
allowed cell and zero elsewhere.  The default seed is uniform on the allowed
support.  Iteration stops only when the maximum observed-marginal residual,
scaled by $\max(1,\text{target})$, meets the declared tolerance.

Missing observations are represented by explicit Boolean masks.  An unobserved
row or column is not scaled and is not interpreted as zero; its fitted value
remains dependent on the seed and other observed constraints.  The result is
therefore marked underdetermined whenever either set of marginals is incomplete.
An explicit observed zero, by contrast, removes all flow from its constrained
row or column.

When all boardings and alightings are observed, their totals must agree.  The
default policy reports unequal totals as infeasible.  The only corrective
policy implemented at this stage is explicit:
\texttt{average\_observed\_totals} rescales each complete marginal vector
proportionally to the arithmetic mean of the two totals.  The original
imbalance, reconciled totals, largest marginal adjustment, and a warning are
retained in the data-quality report; reconciliation is never silent.  A
positive total cannot be reconciled with a zero total.

Complete marginals are checked against the triangular support before
iteration.  In particular, boarding at the final stop and alighting at the
first stop must be zero.  For every destination prefix through stop $k$, its
cumulative alightings cannot exceed boardings at strictly upstream stops:
\[
\sum_{j=1}^{k}a_j
\leq
\sum_{i=1}^{k-1}b_i.
\]
Violations, exhausted positive support, and nonconvergence raise a dedicated
exception carrying the data-quality report.

For a positive fitted boarding marginal, the route-level alighting probability
is
\[
p_{ij}=\frac{x_{ij}}{\sum_{k>i}x_{ik}}.
\]
Zero-flow rows receive zero probabilities.  Results expose immutable matrices,
fitted and reconciled marginals, iterations, absolute and relative residuals,
support size, coverage, and reconciliation diagnostics.

These outputs are explicitly labeled \texttt{leg\_level}, with boarding
semantics \texttt{leg\_boarding\_\allowbreak unclassified}, and are marked incompatible
with journey OD.  A boarding observed on a route may be an initial boarding or
a transfer boarding; route-level counts alone do not identify that distinction.
The IPF result is therefore an audit baseline and initializer, not a procedure
for stitching complete multi-line passenger journeys.

The Phase-3 CPU microbenchmark generated feasible dense triangular marginals.
Routes of 10, 25, 50, and 100 stops converged respectively in 14, 12, 12, and
11 iterations.  Their measured wall times were approximately 0.00059,
0.00078, 0.00131, and 0.00226 seconds, and their dense result matrices occupied
800, 5,000, 20,000, and 80,000 bytes.  These machine-dependent values are not
test thresholds; they show that this diagnostic baseline is negligible
relative to later network journey preprocessing.

